% TOP_PROBLEM: {"id": 3262, "title": "Directed path of length twice the minimum outdegree", "category_id": 3, "category": {"id": 3, "name": "graph_theory", "display_name": "Graph Theory", "description": "Problems involving graphs, networks, and their properties.", "slug": "graph-theory", "order_index": 3, "created_at": "2026-07-31T15:26:25.671Z"}, "status": "open", "problem_number": "OPG-46359", "set_id": 12}
% FIRST_POSTED: 2026-10-02
% ATTEMPT_STATUS: unresolved
% UnsolvedMath ID 3262; source catalogue code OPG-46359
% !TeX program = lualatex
% GPT-6 Astra Ultra research attempt; independently checked partial work, 2026-10-02.
\documentclass[11pt,a4paper]{article}
\usepackage[margin=25mm]{geometry}
\usepackage{fontspec}
\setmainfont{lmroman10-regular.otf}[BoldFont=lmroman10-bold.otf,
 ItalicFont=lmroman10-italic.otf,BoldItalicFont=lmroman10-bolditalic.otf]
\usepackage{amsmath,amssymb,mathtools}
\usepackage{unicode-math}
\setmathfont{latinmodern-math.otf}
\AtBeginDocument{\let\setminus\smallsetminus}
\usepackage{xurl,hyperref}
\hypersetup{colorlinks=true,urlcolor=blue}
\setcounter{secnumdepth}{0}
\setlength{\parindent}{0pt}
\setlength{\parskip}{5pt}
\setlength{\emergencystretch}{3em}
\begin{document}
\section*{Directed path of length twice the minimum outdegree}\label{3262}
{\small UnsolvedMath ID 3262 · OPG-46359 · Graph Theory · OpenGarden}
\subsection{Definitions and mathematical statement}
Let $D=(V,A)$ be a finite nonempty oriented graph, with no loops, no
parallel arcs, and no opposite pair of arcs. Its minimum outdegree is
$\delta^+(D)=\min_{v\in V}|\{w:vw\in A\}|$. A directed path of length
$\ell$ is a sequence of $\ell+1$ distinct vertices
$v_0,v_1,\ldots,v_\ell$ with $v_{i-1}v_i\in A$ for every
$1\le i\le\ell$. Length counts arcs, not vertices.

The conjecture asks whether, for every integer $k\ge1$,
\[
       \delta^+(D)\ge k
       \quad\Longrightarrow\quad
       D\text{ has a directed path of length at least }2k.
\]
One can truncate a longer path, so ``at least $2k$'' is equivalent to
existence of a path of length exactly $2k$. No strong-connectivity
hypothesis or lower bound on indegrees is present. The word ``path''
in the target is essential despite the historical source URL containing
``cycle''.

A simple necessary order condition is automatic. Since an oriented
$n$-vertex graph has at most $n(n-1)/2$ arcs, minimum outdegree $k$
implies $n\ge2k+1$, leaving enough vertices for the requested path.
A regular tournament on $2k+1$ vertices shows that a larger universal
path length cannot be demanded from this outdegree hypothesis alone.
The conjecture requires distinct vertices throughout; a directed walk
with repeated vertices would not answer it.
\subsection{Short English statement}
Must an oriented graph of minimum outdegree k contain a directed path with 2k arcs and 2k+1 distinct vertices?
\subsection{Research attempt: an endpoint cycle and its exit}
\paragraph{Scope and literature check.}
GPT-6 Astra Ultra, 2 October 2026. Path length counts arcs. There is no
indegree assumption, and an oriented graph has no opposite arcs.
Cheng and Keevash [S3] prove a path-length lower bound of
$1.5\delta^+$ in this class. Their counterexamples to a stronger
higher-girth variant do not disprove the present twice-outdegree
conjecture. The elementary argument below is weaker in general,
but it isolates a concrete cycle-extension mechanism, settles
$k=1,2$, and gives the full conclusion for tournaments.

\paragraph{1. A long cycle from a longest path.}
Choose a longest directed path $P=(v_0,\ldots,v_\ell)$. Every
outneighbour of its last vertex lies on $P$, since otherwise the
path extends. Neither $v_\ell$ nor $v_{\ell-1}$ is an outneighbour:
the first exclusion is looplessness and the second is the prohibition
of opposite arcs. There are at least $k$ outneighbours among indices
$0,\ldots,\ell-2$. If $j$ is the smallest such index, counting the
available positions gives $j\le\ell-k-1$. Hence
\[
       v_j,v_{j+1},\ldots,v_\ell,v_j
\]
is a directed cycle with at least $k+2$ vertices. This proof also
shows why dropping the orientation hypothesis changes the bound:
a return arc to the predecessor would no longer be excluded.
Deleting one arc of the cycle gives a path with at least $k+1$ arcs.

\paragraph{2. A cycle of at most $2k$ vertices has an exit.}
Let $C$ be any directed cycle with $c\le2k$ vertices. If every
outgoing arc from a vertex of $C$ stayed inside $V(C)$, its induced
oriented subgraph would have at least $ck$ arcs. But an oriented
$c$-vertex graph has at most $c(c-1)/2<ck$ arcs, a contradiction.
There is therefore an arc $u\to w$ with $u\in V(C)$ and
$w\notin V(C)$.

Start at the successor of $u$ around $C$, follow the cycle through
all its vertices until reaching $u$, and then take $u\to w$.
This is a simple directed path with exactly $c$ arcs. In particular,
any directed cycle of length at least $2k$ yields the desired
$2k$-arc path: use its exit when its length is exactly $2k$, and
break the cycle when its length is larger.

Combining the two lemmas proves the full statement for $k=1,2$.
For $k=1$ the first cycle has length at least three and already
contains a two-arc path. For $k=2$ its length is at least four; a
four-cycle has an exit, and a longer cycle directly supplies four
arcs. For every $k\ge2$ the same reasoning gives a path with at least
$k+2$ arcs: either the cycle has length at most $2k$ and has an exit,
or its length is at least $2k+1$ and supplies $2k\ge k+2$ arcs.

\paragraph{3. The tournament subclass.}
Every tournament has a Hamilton path, by induction on its order.
Given a directed path through all vertices except $x$, put $x$
before the first path vertex it dominates. If there is a predecessor
of that vertex, the predecessor dominates $x$ by the choice of the
first position, so both new path arcs exist. If there is no dominated
vertex, append $x$ at the end. This inserts $x$ in all cases.

A tournament with minimum outdegree at least $k$ has order at least
$2k+1$, by counting its $n(n-1)/2$ arcs. Its Hamilton path consequently
has at least $2k$ arcs. This handles all $k$ in that subclass without
an indegree hypothesis. In a general oriented graph a missing
adjacency invalidates precisely the insertion step above.

\paragraph{Verification and first remaining gap.}
The script [S9] exhausts all labelled oriented graphs through
four vertices for the endpoint-cycle lemma, and checks tournament
insertion through five vertices. These bounded checks supplement
the general proofs. For $k>2$, the cycle obtained at the endpoint
may have fewer than $2k$ vertices. One exit then supplies only its
own length in path arcs, with no guarantee of a second compatible
extension. Controlling further extensions is the first gap.

\paragraph{Final status.}
UNRESOLVED generally. The cases $k=1,2$, every tournament, and the
stated elementary lower bound are proved, without claiming to
improve the stronger general bound in [S3].
\subsection{Sources}
\begin{itemize}
\item[S1] ULAM.AI and the UnsolvedMath Contributors, supplied problems.json, record 3262 (OPG-46359), OpenGarden collection. Catalogue identity and source transcription; CC BY 4.0.
\url{https://huggingface.co/datasets/ulamai/UnsolvedMath}.
\item[S2] Open Problem Garden, Directed path of length twice the minimum outdegree. Original problem and discussion; the mathematical formulation below is expanded from the supplied source record.
\url{http://www.openproblemgarden.org/op/directed_cycle_of_length_twice_the_minimum_outdegree}.
\item[S3] Yangyang Cheng and Peter Keevash, \emph{On the length of directed paths in digraphs}, SIAM Journal on Discrete Mathematics (2025). \url{https://arxiv.org/abs/2402.16776}; \url{https://doi.org/10.1137/24M1648375}.
\item[S9] GPT-6 Astra Ultra, exact finite checks accompanying this attempt (2 October 2026). Repository script, Python standard library only. \texttt{scripts/check\_attempt\_batch03\_digraphs\_2026\_10\_02.py}.
\end{itemize}
\end{document}
